%% Question bank chapter file - generated by bank_tools.py build (do not edit while the bank is still being built:
%% it is regenerated from the per-test data in Question Bank\data). Every question is one bankq block:
%%   \begin{bankq}[id=..., source={...}, paper=.., qno=.., type=..., subtopic={...}, answer={...}, answer_src=key|solved,
%%                 working={...}, difficulty=1-10, diff_parts={K S I C F}, diff_note={...},
%%                 group=<passage id or empty>, also_in={duplicate copies in other tests}]
%% difficulty = min(10, K+S+I+C+F-1): K concepts 1-3, S steps 1-4, I insight 0-2, C computation 0-2, F format 0-1
%%   ...question LaTeX, options, figures (\bfig / \bfigi -> img/)...
%%   \end{bankq}
%% Passages shared by several questions are bankpassage blocks; their questions carry group=<passage id>.
\documentclass[10pt]{article}
\usepackage[a4paper,margin=1.8cm]{geometry}
\usepackage[T1]{fontenc}
\usepackage{amsmath,amssymb,mathtools}
\usepackage{graphicx}
\usepackage[export]{adjustbox}
\usepackage{xcolor,needspace,keyval}
\usepackage{enumitem}
\usepackage{booktabs,array,tabularx}
\usepackage[hidelinks]{hyperref}
\setlength{\parindent}{0pt}
\setlength{\parskip}{3pt}
\makeatletter
\define@key{bq}{id}{\def\bq@id{#1}}\define@key{bq}{source}{\def\bq@src{#1}}\define@key{bq}{paper}{\def\bq@paper{#1}}
\define@key{bq}{qno}{\def\bq@qno{#1}}\define@key{bq}{type}{\def\bq@type{#1}}\define@key{bq}{subtopic}{\def\bq@sub{#1}}
\define@key{bq}{answer}{\def\bq@ans{#1}}\define@key{bq}{answer_src}{\def\bq@asrc{#1}}\define@key{bq}{working}{\def\bq@work{#1}}
\define@key{bq}{group}{\def\bq@group{#1}}\define@key{bq}{also_in}{\def\bq@also{#1}}
\define@key{bq}{difficulty}{\def\bq@diff{#1}}\define@key{bq}{diff_parts}{\def\bq@dparts{#1}}\define@key{bq}{diff_note}{\def\bq@dnote{#1}}
\newenvironment{bankq}[1][]{\def\bq@work{}\def\bq@also{}\def\bq@group{}\def\bq@diff{?}\def\bq@dparts{}\def\bq@dnote{}\setkeys{bq}{#1}%
  \par\needspace{4\baselineskip}\medskip\noindent{\small\textbf{\bq@id}\hfill\color{gray}[\bq@src, Paper \bq@paper, Q\bq@qno]
  \bq@type\ $\cdot$ \bq@sub\ $\cdot$ \textbf{D\bq@diff}}\par\nopagebreak\relax}
  {\par{\small\color{teal}Answer: \textbf{\bq@ans}\ifx\bq@work\@empty\else\ (solved: \bq@work)\fi
  \ifx\bq@also\@empty\else\par\color{gray}Also in: \bq@also\fi}\par}
\define@key{bp}{id}{\def\bp@id{#1}}\define@key{bp}{source}{\def\bp@src{#1}}\define@key{bp}{questions}{\def\bp@qs{#1}}
\newenvironment{bankpassage}[1][]{\setkeys{bp}{#1}\par\needspace{4\baselineskip}\bigskip\noindent
  {\small\textbf{Passage \bp@id}\hfill\color{gray}[\bp@src; questions \bp@qs]}\par\nopagebreak\relax\itshape}{\par}
\makeatother
\newcommand{\bfig}[2]{\par\smallskip{\centering\includegraphics[width=#1pt,max width=0.62\linewidth,max height=0.32\textheight]{img/#2}\par}\smallskip}
\newcommand{\bfigi}[2]{\includegraphics[width=#1pt,max width=\linewidth,max height=0.32\textheight,valign=c]{img/#2}}
\newcommand{\bmissing}[1]{\textit{[figure missing in the source conversion: \detokenize{#1}]}}
\newenvironment{qoptions}{\begin{itemize}[leftmargin=1.8em,itemsep=1pt,topsep=2pt,label={}]}{\end{itemize}}
\newcommand{\opt}[2]{\item[(#1)] #2}
\providecommand{\fwstart}{}\providecommand{\fwend}{}\providecommand{\qfullsplit}{}
\newcolumntype{L}[1]{>{\raggedright\arraybackslash}p{#1}}
\newcolumntype{C}[1]{>{\centering\arraybackslash}p{#1}}
\newcolumntype{R}[1]{>{\raggedleft\arraybackslash}p{#1}}
\providecommand{\sectionhdr}[1]{}\providecommand{\subsectionhdr}[1]{}


\begin{document}
{\Large\bfseries Magnetism --- Question Bank}\par\smallskip
{\small Physics chapter P11. Questions from all converted test papers, grouped by source series. Answers: printed key where it matches the paper, otherwise solved (answer\_src=solved, with working). Exact repeats across papers appear once, with also\_in listing the other copies.}\par\medskip\hrule

\section*{Score-I papers}
\begin{bankq}[id=SI12-P1-PHY-S1-Q03, source={Score-I - Test 12 (09 Mar 2022)}, paper=1, qno=3, type=multi, subtopic={Force on charges}, answer={A,C,D}, answer_src=key, difficulty=4, diff_parts={K2 S2 I0 C0 F1}, diff_note={Charged particle under uniform force + projectile kinematics: parabola, time of flight, reduced range}]
A particle of mass $m$ and charge $-q$ has been projected from ground as shown in the figure below. Mark out the correct statement(s):
\bfig{92.0}{SI12_p1_phy_s1_q03_1.png}
\begin{qoptions}
\opt{A}{The path of motion of the particle is parabolic}
\opt{B}{The path of motion of the particle is a straight line}
\opt{C}{Time of flight of particle is $\dfrac{2u\sin\theta}{g}$}
\opt{D}{Range of motion of the particle will be less than $\dfrac{u^2\sin2\theta}{g}$}
\end{qoptions}
\end{bankq}
\begin{bankq}[id=SI12-P1-PHY-S2-Q01, source={Score-I - Test 12 (09 Mar 2022)}, paper=1, qno=1, type=numerical, subtopic={Force on wires}, answer={1.00}, answer_src=key, difficulty=4, diff_parts={K2 S2 I0 C0 F1}, diff_note={Field of long wire + force on the two parallel loop sides (opposite sides cancel)}]
An infinitely long straight wire carries a current $i_1=1$ amp. A rectangular current loop placed in the same plane of shown dimensions carries a current $i_2=1$ amp anticlockwise. Magnitude of net force on the loop is found to be $\beta\dfrac{\mu_0}{4\pi}$. Find value of $\beta$.
\bfig{128.9}{SI12_p1_phy_s2_q01_1.png}
\end{bankq}
\begin{bankq}[id=SI12-P2-PHY-S2-Q03, source={Score-I - Test 12 (09 Mar 2022)}, paper=2, qno=3, type=numerical, subtopic={Force on charges}, answer={2.50}, answer_src=key, difficulty=5, diff_parts={K2 S2 I1 C0 F1}, diff_note={Circle through origin: R^2=a^2+(R-d)^2 gives R=25/8, p=qBR}]
A particle of charge $+q$ enters a region of uniform magnetic field $B$ (directed into the plane of paper) as shown in the figure. Particle is deflected by distance $d$ along $Y$ axis after travelling a distance of $a$ along $X$ axis. Find magnitude of linear momentum of particle in Newton-sec. (Given: $a=3$ m, $d=4$ m, $qB=0.8$ C-Tesla)
\bfig{150.4}{SI12_p2_phy_s2_q03_1.png}
\end{bankq}
\begin{bankq}[id=SI12-P2-PHY-S2-Q04, source={Score-I - Test 12 (09 Mar 2022)}, paper=2, qno=4, type=numerical, subtopic={Force on charges}, answer={6.02 to 6.03}, answer_src=key, difficulty=8, diff_parts={K2 S3 I2 C1 F1}, diff_note={Equal radii circles tangent at origin; angular speeds ratio 24 so meeting when A has turned 2pi*24/25}]
A neutral atom of atomic mass number 100 which is stationary at the origin in gravity free space emits an $\alpha$-particle ($A$) in $z$-direction. The product ion is $P$. A uniform magnetic field exists in the $x$-direction. Disregard the electro magnetic interaction between $A$ and $P$. If the angle of rotation of $A$ after which $A$ and $P$ will meet for the first time is $\theta$ radians, what is the value of $\theta$?
\end{bankq}
\end{document}
